\section{How This Book Was Made}
\label{app:how-made}

This appendix documents the engineering behind the book itself: the CrewAI
agent pipeline that produced the content and the LaTeX toolchain that compiled
it. It is included to demonstrate the architecture, not merely describe it
\cite{team2026}.

\subsection{The agent pipeline}

The book is generated by a sequential CrewAI crew. Each agent maps to a job
title in a small publishing house, and the output of one stage flows as context
into the next. The three language-model stages route every call through a
central API gatekeeper that enforces rate limits, retries, and cost logging.

\begin{figure}[htbp]
\centering
\begin{tikzpicture}[
  node distance=7mm and 16mm,
  stage/.style={draw, rounded corners, align=center, font=\small,
                minimum width=42mm, minimum height=10mm, fill=blue!5},
  side/.style={draw, dashed, rounded corners, align=center, font=\footnotesize,
               minimum width=30mm, minimum height=9mm, fill=gray!8},
  arr/.style={-{Stealth[length=2.4mm]}, thick},
  dep/.style={-{Stealth[length=2mm]}, dashed, gray}
]
\node[stage] (rd) {1. Research Director};
\node[stage, below=of rd] (oa) {2. Outline Architect};
\node[stage, below=of oa] (cw) {3. Chapter Writer};
\node[stage, below=of cw] (asm) {4. LaTeX Assembly};
\node[stage, below=of asm] (rev) {5. Review \& Page Check};
\node[stage, below=of rev] (cmp) {6. Compile PDF};

\draw[arr] (rd) -- (oa);
\draw[arr] (oa) -- (cw);
\draw[arr] (cw) -- (asm);
\draw[arr] (asm) -- (rev);
\draw[arr] (rev) -- (cmp);

\node[side, right=of oa] (gk) {API Gatekeeper\\(limits, retries, cost)};
\draw[dep] (rd.east) -- (gk.west);
\draw[dep] (oa.east) -- (gk.west);
\draw[dep] (cw.east) -- (gk.west);

\node[side, right=of cmp] (pdf) {\texttt{book.pdf}};
\draw[arr] (cmp.east) -- (pdf.west);
\end{tikzpicture}
\caption{CrewAI book-generation pipeline. Solid arrows show the sequential
hand-off of artifacts; dashed arrows show language-model calls routed through
the gatekeeper.}
\label{fig:pipeline}
\end{figure}

\subsection{From IPC orchestration to a CrewAI crew}

An earlier exercise orchestrated agents directly with inter-process
communication, wiring every queue, retry, and hand-off by hand. That approach
is flexible and fast for a proof of concept, but it pushes all coordination
onto the developer. CrewAI moves that boilerplate into a framework: roles,
tasks, and a process type replace custom message plumbing, so the team can
focus on prompts, schemas, and quality gates. The trade-off is less low-level
control in exchange for clearer structure and easier hand-over.

\subsection{Observability and cost}

Because a production agent is a system rather than a single prompt, every
language-model call is logged as a JSON line with token counts and an estimated
cost, and the gatekeeper records any request that exceeds a configured budget.
This makes a run reproducible and auditable: the same configuration yields the
same book, and the cost of a live run can be reported from the logs rather than
guessed.
